% ---------------------------------------------------------------------------
% Identity tiles for every large length, from two conditions at one point.
% AI-DRAFTED on 4 October 2026 (assistance record, §10).  The mathematics is
% machine-verified (verification/tile_controllability.py, tile_wronskian.py,
% tile_rank_exact.py); the prose is to be rewritten by the author.
% ---------------------------------------------------------------------------
\subsection*{Existence of identity tiles without finding them}
\label{sec:tilesallk}

The controllability question above can be answered. The answer rests on two
facts about a single, explicitly chosen tile --- the tile whose interior
\emph{is} the dock --- and the rest is an argument about compact subgroups.
It yields $Z(P(n,k))=M(P(n,k))=2k+2$ for all sufficiently large $n$ at every
$k$ for which the two facts are verified, and they are verified exactly for
$k\le\TileAllK$; in particular at $k=6,\dots,\TileAllK$, where no identity tile has
ever been found. The price is that the threshold $L(k)$ is not effective.

Throughout, a \emph{dock} is a constant weight pattern $(a_0,b_0,c_0,d_0,e_0)$
with $b_0c_0e_0\ne0$, normalised to $b_0=1$; the \emph{dock symbol} is the
rotation-invariant symbol of \S\ref{sec:cyclo} with these weights,
\[
  F_{\mathrm{dock}}(s)=(a_0+s)\bigl(d_0+e_0\,t_k(s)\bigr)-c_0^{2},
  \qquad t_k(s)=2T_k(s/2),
\]
and $S_0$ is the common step matrix of the constant weight sequence, so that a
tile whose interior equals the dock has product $S_0^{\ell}$.

\begin{lemma}[The invariant form, exactly]\label{lem:wronskian}
Fix a dock. For the state $Z_0=(x_{-k-1},\dots,x_k)$ let $B=B(\mathrm{dock})$ be
the skew form
\[
  Z^{\!\top}BZ' \;=\; \sum_{\substack{\text{edges }pq\\ p\le -1<q}} A_{pq}\,(z_pz'_q-z'_pz_q),
\]
the flux of two kernel sequences across the cut between positions $-1$ and $0$,
with inner coordinates eliminated by $y_m=-(x_{m-1}+a_0x_m+x_{m+1})/c_0$. Then
$P^{\!\top}BP=B$ for every tile product $P$, and
\[
  \det B \;=\; e_0^{\,2k}/c_0^{\,4k}\;\ne\;0 .
\]
Hence every tile product lies in $Sp(B)\cong Sp(2k+2,\mathbb R)$.
\end{lemma}

\begin{proof}
For two kernel sequences $z,z'$ of the strip operator, the flux across the cut
between $j$ and $j+1$ changes, when the cut moves one position, by
$\sum_p\bigl(z_p(Az')_p-z'_p(Az)_p\bigr)$ over the vertices that cross, which is
$0$. So the flux is the same across every cut. Across the cut $(-1,0)$ it
involves the outer edge $u_{-1}u_0$ and the $k$ inner edges $v_mv_{m+k}$ with
$-k\le m\le-1$, whose inner coordinates are linear in $x_{-k-1},\dots,x_k$:
this is the form $B$, and it depends only on the weights at positions
$-k-1,\dots,k$. Transporting the state from the cut $(-1,0)$ to the cut
$(\ell-1,\ell)$ by the tile product preserves the flux, and at both cuts the
weights in range are dock weights (the tile's own dock at one end, the
neighbouring tile's at the other), so both forms are $B$: $P^{\!\top}BP=B$.
The determinant is computed in exact arithmetic for the docks used below
(\texttt{verification/tile\_wronskian.py}); the closed form was checked on
random rational docks for $k\le8$ and is stated as what it is.
\end{proof}

\begin{lemma}[Elliptic docks exist for every $k$]\label{lem:elliptic}
Let $\varepsilon_k=1$ for $k$ even and $\varepsilon_k=(-1)^{(k-1)/2}$ for $k$ odd.
For $c_0>0$ small enough the dock $(0,1,c_0,0,\varepsilon_k)$ has
$F_{\mathrm{dock}}(s)=\varepsilon_k\,s\,t_k(s)-c_0^{2}$ with $k+1$ distinct real
roots in $(-2,2)$; hence $S_0$ is diagonalisable with all eigenvalues on the
unit circle, and $\{S_0^m:m\ge0\}$ has compact closure, a group.
\end{lemma}

\begin{proof}
$g(s)=s\,t_k(s)$ has the $k$ simple zeros $2\cos\frac{(2j-1)\pi}{2k}$ of $t_k$
in $(-2,2)$ and the zero $s=0$. For even $k$ these are $k+1$ distinct simple
zeros, $g(\pm2)=\pm4\ne0$, and simple zeros persist under the small
perturbation $-c_0^2/\varepsilon_k$. For odd $k$, $t_k$ is odd, so $s=0$ is a
double zero of $g$ with $g(s)=t_k'(0)s^2+O(s^4)$ and $t_k'(0)=k(-1)^{(k-1)/2}$;
the equation $\varepsilon_kg(s)=c_0^2$ then has two real roots $\pm c_0/\sqrt{k}+O(c_0^3)$
near $0$ because $\varepsilon_kt_k'(0)=k>0$, and the other $k-1$ zeros are simple
and persist. The eigenvalues of $S_0$ are the $z$ with $F_{\mathrm{dock}}(z+z^{-1})=0$;
real roots $s\in(-2,2)$ give $z=e^{\pm i\theta}$ with $s=2\cos\theta$, distinct
for distinct $s$. A diagonalisable matrix with unimodular eigenvalues generates
a relatively compact cyclic group, whose closure is a compact abelian group.
\end{proof}

\begin{theorem}[Identity tiles for every large length]\label{thm:tilesallk}
Fix $k\ge2$ and a dock satisfying Lemma~\ref{lem:elliptic}. Suppose that for
some length $\ell_0$ the differential of the tile map $w\mapsto P_{\ell_0}(w)$ has
rank $(k+1)(2k+3)=\dim Sp(2k+2)$ at some point. Then there is $L$ such that an
identity tile of every length $\ell\ge L$ exists. Consequently
$Z(P(n,k))=M(P(n,k))=2k+2$ for all $n\ge L$.
\end{theorem}

\begin{proof}
Write $G=Sp(B)$, a connected Lie group by Lemma~\ref{lem:wronskian}, and fix a
left-invariant metric. The set where the differential of $P_{\ell_0}$ has rank
$<\dim G$ is a proper algebraic subset of the weight space (the rank is
maximal on an open dense set), so there are full-rank points $w^*$ arbitrarily
close to the dock point $w_0$ (interior $=$ dock), and at such points all edge
weights are nonzero. At a full-rank point $P_{\ell_0}$ is a submersion onto $G$,
hence open: $S_{\ell_0}$ contains a ball $B_\varepsilon(P^*)$ around
$P^*=P_{\ell_0}(w^*)$. Since $P^*$ is close to $S_0^{\ell_0}$ and
$S_0^{\ell_0}$ is diagonalisable with distinct unimodular eigenvalues --- an
open condition in $Sp$, because unimodular eigenvalues of a symplectic matrix
can leave the circle only by colliding --- we may take $P^*$ elliptic: the
closure $C$ of $\{P^{*m}:m\in\mathbb Z\}$ is compact.

Let $\kappa$ bound the Lipschitz constant of conjugation by elements of $C$ on
a neighbourhood of the identity, and put $B'=B_{\varepsilon/\kappa}(I)$. For
$v_1,\dots,v_M\in B'$ set $u_i=P^{*\,M-i}v_iP^{*\,-(M-i)}\in B_\varepsilon(I)$,
so that $P^*u_i\in B_\varepsilon(P^*)\subseteq S_{\ell_0}$ and, telescoping,
\[
  (P^*u_1)(P^*u_2)\cdots(P^*u_M)=P^{*M}\,v_1v_2\cdots v_M .
\]
The sets $B'^M$ increase with $M$ and exhaust an open subgroup of $G$, which is
$G$ itself; the set $K=\{P^{*-M}S_0^{-r}: M\ge0,\ 2k+2\le r<2k+2+\ell_0\}$ is
contained in $C\cdot\{S_0^{-r}\}$, a compact set, so $K\subseteq B'^{M_0}$ for
some $M_0$. Given $N\ge M_0\ell_0+2k+2$ write $N=M\ell_0+r$ with $M\ge M_0$ and
$2k+2\le r<2k+2+\ell_0$, choose $v_i\in B'$ with $v_1\cdots v_M=P^{*-M}S_0^{-r}$,
and concatenate the $M$ tiles $P^*u_i$ with a pure-dock tile of length $r$
(product $S_0^{r}$): by Lemma~\ref{lem:concat} the result is a tile of length
$N$ with product $P^{*M}\,v_1\cdots v_M\,S_0^{r}=I$. Theorem~\ref{thm:tiling}
with Theorem~\ref{thm:red} gives the nullity $2k+2$ and hence
$Z=M=2k+2$ for every $n\ge L:=M_0\ell_0+2k+2$.
\end{proof}

\begin{remark}[Three directions per position]\label{rem:threeperpos}
The hypothesis is checkable, and the length at which it holds is predictable.
Rescaling the coordinate at $u_i$ and at $v_i$ by nonzero constants is a
diagonal similarity $A\mapsto DAD$, which preserves the pattern, the nullity
and the monodromy up to conjugation, and moves two of the five weights at
position $i$. So only three weights per position act on the tile product, and
the generic rank of the differential is
\[
  \operatorname{rank}dP_\ell \;=\; 3(\ell-2k-2)+(k+1) \;=\; 3\ell-5k-5
\]
until it saturates at $\dim Sp=(k+1)(2k+3)$, a law we observe exactly at every
$(k,\ell)$ computed. The first submersive length is therefore
$\ell_F(k)=\lceil(k+1)(2k+8)/3\rceil$: $19,27,36,47,\dots$ for $k=3,4,5,6,\dots$,
and the identity tiles of Theorem~\ref{thm:k3all} first appeared at $21$ and
$29$ --- two above $\ell_F$ --- which also explains why the count
$5(\ell-2k-2)\ge(2k+2)^2$ used earlier overestimated the parameters available.
At the dock point itself the rank is $\dim Sp-(k-2)$: a constant perturbation of
a constant background moves the $k+1$ eigenphases of $S_0$ only through the
three parameters of the symbol, the obstruction of \S\ref{sec:cyclo} in
infinitesimal form.
\end{remark}

\begin{proposition}[The hypothesis, verified exactly]\label{prop:tilesallk}
For $2\le k\le\TileAllK$, with the integer docks of
Table~\ref{tab:tilesallk}: the dock symbol has $k+1$ distinct real roots in
$(-2,2)$ (Sturm's theorem over $\mathbb Q$); $\det B(\mathrm{dock})\ne0$ (exact
rational arithmetic); and at the integer point $w^*$ of the table the
differential of $P_{\ell_0}$, computed by exact forward-mode differentiation
over $\mathbb F_p$, $p=33554393$, has rank $(k+1)(2k+3)$ --- a lower bound for
the rank over $\mathbb Q$. Hence Theorem~\ref{thm:tilesallk} applies, and
\[
  Z(P(n,k))=M(P(n,k))=2k+2\qquad\text{for all sufficiently large }n,\ \ 2\le k\le\TileAllK .
\]
\end{proposition}

\begin{table}[ht]\centering\small
% GENERATED by verification/verify_all.py -- do not edit by hand.
\begin{tabular}{r l r r r}\hline
$k$ & dock $(a_0,c_0,d_0,e_0)$, $b_0=1$ & $\ell_0$ & rank $=\dim Sp$ & $\det B$\\\hline
$2$ & $(0,1,0,1)$ & $15$ & $21=21$ & $1$\\
$3$ & $(1,1,1,1)$ & $21$ & $36=36$ & $1$\\
$4$ & $(0,1,1,1)$ & $29$ & $55=55$ & $1$\\
$5$ & $(0,1,0,1)$ & $39$ & $78=78$ & $1$\\
$6$ & $(0,1,-1,2)$ & $49$ & $105=105$ & $4096$\\
$7$ & $(0,1,0,-1)$ & $61$ & $136=136$ & $1$\\
$8$ & $(-1,1,0,2)$ & $75$ & $171=171$ & $65536$\\
$9$ & $(0,1,0,2)$ & $89$ & $210=210$ & $262144$\\
\hline\end{tabular}

\caption{The docks and lengths at which the hypothesis of
Theorem~\ref{thm:tilesallk} is certified. Generated by
\texttt{verification/verify\_all.py}.}
\label{tab:tilesallk}
\end{table}

\begin{remark}[What is new, and what is not]
For $k=3,4,5$ the statement was already known with effective thresholds
(Theorems~\ref{thm:k3all}, \ref{thm:k5all}) by finding and certifying the
tiles; for $k\ge6$ no tile has been found, and Proposition~\ref{prop:tilesallk}
is the first statement of any kind about $Z(P(n,k))$ for all large $n$ at those
$k$. The bound $L(k)$ is ineffective because $M_0$ comes from compactness.
The rank hypothesis is proved for every $k$ in \S\ref{sec:tilesrank}: the
$k-2$ directions missing at the dock point are supplied at second order.
\end{remark}
